%!TEX root = ../main.tex
\chapter{The Interpreter}

Sameer had worked as a court interpreter for eleven years, and in that
time he had learned that the hardest sentences to translate were never
the technical ones. Legal terminology had a fixed and knowable shape;
you learned the vocabulary once and it stayed learned. What broke a
translation, what made a witness's testimony come out flattened and
wrong on the other side of it, was always something smaller: the
particular way a mother said her son's name when she still believed he
was innocent, the half-second hesitation before a lie.

Today's case was a property dispute, ordinarily the most tedious kind
of work, full of boundary lines and inheritance clauses. But the
plaintiff, an elderly man named Younes, kept drifting from the
questions he was asked into a story about the house itself — who had
built it, which room his father had died in, the fig tree his
grandmother had planted the year the family arrived.

``Mr. Younes,'' the lawyer said, not unkindly, ``can you confirm the
date the property was transferred to your name?''

Younes answered in a long paragraph about the fig tree.

Sameer translated the date, because that was, strictly, the answer
buried inside the paragraph. But he found himself adding, quietly,
after the official translation: ``He says the tree is still there. He
says that's how he knows the house is still his, whatever the papers
say.''

The lawyer frowned at the addition. The judge did not. She looked at
Younes for a long moment, and then at Sameer, and said, ``Let the
witness continue.''

It was not, strictly, his job to decide what mattered. But eleven
years in this room had taught him that the truth of a testimony rarely
lived in the part anyone had thought to ask about.
